\documentclass[conference]{IEEEtran}
\usepackage{graphicx}
\usepackage{float}
\usepackage{placeins}
\usepackage{amsmath}
\affiliation{\institution{Example University}\country{USA}}
\email{author@example.edu}

\graphicspath{{./}}

\begin{document}
\title{Graph Neural Networks: A Lightweight Survey}
\author{W. Hamilton}
\maketitle
\begin{abstract}
We summarize message passing networks and common training objectives.
\end{abstract}
\begin{IEEEkeywords}
graph neural networks, message passing
\end{IEEEkeywords}
\IEEEoverridecommandlockouts
\section{Introduction}
Graph neural models exchange information over local neighborhoods.

\FloatBarrier
\section{Message Passing}
At layer $l+1$:
\begin{equation}
h_v^{(l+1)}=\phi\left(h_v^{(l)},\square_{u\in\mathcal{N}(v)}\psi(h_v^{(l)},h_u^{(l)})\right).
\end{equation}

\FloatBarrier
\section{Conclusion}
Model design is a balance between expressivity and efficiency.

\FloatBarrier
\bibliographystyle{IEEEtran}
\bibliography{references}
\end{document}
